\documentclass[landscape,a4paper, english, 10pt]{scrartcl}
\usepackage[utf8]{inputenc}
\usepackage[english]{babel}
\usepackage[T1]{fontenc}
\usepackage{tikz} 
\usepackage{translator}
\usepackage{fancyhdr}
\usepackage{fix-cm}
\usepackage[landscape, headheight = 2cm, margin=.5cm, top = 3.2cm, nofoot]{geometry}
\usepackage[dvipsnames]{xcolor}

\usetikzlibrary{calc, calendar}

\renewcommand*\familydefault{\sfdefault}

% 2026年設定
\def\year{2026}

% 【修正ポイント1】ラベル位置の微調整
% yshiftを 0em にし、anchorを west にすることで日付ボックスのすぐ右側に配置します
\newcommand{\termin}[2]{
  \node [anchor=west, text width= 3.4cm, xshift=3.2em] at
    (cal-#1.west) {\tiny{#2}};
}

% ヘッダー
\renewcommand{\headrulewidth}{0.0pt}
\setlength{\headheight}{10ex}
\chead{
  \fontsize{60}{70}\selectfont\textbf{\year}
  \Large\textbf{Eduop Chiba}\hfill
}
%%%%%%%%%%%%%%%%%%%%%%%
% パッケージの読み込み群に追加してください
\usepackage{luatexja}
% 予定を追加するマクロ（第1引数: YYYYMMDD, 第2引数: 予定のテキスト）
% --- 予定を追加するマクロ（LuaLaTeX対応版） ---
\makeatletter
\newcommand{\addschedule}[2]{%
  \@addschedule#1\relax{#2}%
}
\def\@addschedule#1#2#3#4#5#6#7#8\relax#9{%
  % CJK環境は不要です。そのまま #9 を出力します。
  \node [anchor=west, text width= 3.4cm, xshift=3.2em, text=black] at
    (cal-#1#2#3#4-#5#6-#7#8.west) {\tiny\textbf{#9}};
}
\makeatother
%%%%%%%%%%%%%%%%%%%%%%%%
\begin{document}
\pagestyle{fancy}
\begin{center}
\begin{tikzpicture}[every day/.style={anchor = north}]
\calendar[
  dates=\year-01-01 to \year-06-30,
  name=cal,
  day yshift = 3em,
  day code={
    \node[name=\pgfcalendarsuggestedname,every day,shape=rectangle,
    minimum height= .53cm, text width = 4.4cm, draw = gray!30]{\tikzdaytext};
    \draw (-1.8cm, -.1ex) node[anchor = west]{\footnotesize%
      \pgfcalendarweekdayshortname{\pgfcalendarcurrentweekday}};
  },
  execute before day scope={
    \ifdate{day of month=1}{
      \pgftransformxshift{4.8cm}
      \draw (0,0)node [shape=rectangle, minimum height= .53cm,
        text width = 4.4cm, fill = NavyBlue!80, text= white, draw = NavyBlue!80, text centered]
        {\textbf{\pgfcalendarmonthname{\pgfcalendarcurrentmonth}}};
    }{}
    \ifdate{workday}{\tikzset{every day/.style={fill=white}}}{}
     \ifdate{between=2026-01-01 and 2026-01-03}{%
       \tikzset{every day/.style={fill=gray!30}}}{}
   \ifdate{between=2026-03-23 and 2026-04-24}{%
       \tikzset{every day/.style={fill=gray!30}}}{}
    \ifdate{Saturday}{\tikzset{every day/.style={fill=blue!5}}}{}
    \ifdate{Sunday}{\tikzset{every day/.style={fill=red!10}}}{}
    % --- 2026年 祝日判定の正確な修正 ---
    \ifdate{equals=01-01}{\tikzset{every day/.style={fill=red!10}}}{}
    \ifdate{equals=01-12}{\tikzset{every day/.style={fill=red!10}}}{} % 成人の日(1/12)
    \ifdate{equals=02-11}{\tikzset{every day/.style={fill=red!10}}}{}
    \ifdate{equals=02-23}{\tikzset{every day/.style={fill=red!10}}}{}
    \ifdate{equals=03-20}{\tikzset{every day/.style={fill=red!10}}}{}
    \ifdate{equals=04-29}{\tikzset{every day/.style={fill=red!10}}}{}
    \ifdate{equals=05-03}{\tikzset{every day/.style={fill=red!10}}}{}
    \ifdate{equals=05-04}{\tikzset{every day/.style={fill=red!10}}}{}
    \ifdate{equals=05-05}{\tikzset{every day/.style={fill=red!10}}}{}
    \ifdate{equals=05-06}{\tikzset{every day/.style={fill=red!10}}}{}
  },
  execute at begin day scope={\pgftransformyshift{-.53*\pgfcalendarcurrentday cm}}
];
% 【修正ポイント2】PDFでズレていた日付を2026年の正式なものに固定
\termin{\year-01-01}{New Year's Day}
\termin{\year-01-12}{Coming of Age Day}
\termin{\year-02-11}{National Foundation Day}
\termin{\year-02-23}{Emperor's Birthday}
\termin{\year-03-20}{Vernal Equinox Day}
\termin{\year-04-29}{Showa Day}
\termin{\year-05-03}{Constitution Day}
\termin{\year-05-04}{Greenery Day}
\termin{\year-05-05}{Children's Day}
\termin{\year-05-06}{Sub. Holiday}
%%%%%%%%%%%%%%%%%%%%%%%%%%%%%%%%%%%%
%予定を追加
\addschedule{20260319}{終業式}
\addschedule{20260427}{始業式}
\end{tikzpicture}
%%%%%%%%%%%%%%%%%%
\pagebreak
%%%%%%%%%%%%%%%%%
\begin{tikzpicture}[every day/.style={anchor = north}]
\calendar[
  dates=\year-07-01 to \year-12-31,
  name=cal,
  day yshift = 3em,
  day code={
    \node[name=\pgfcalendarsuggestedname,every day,shape=rectangle, 
      minimum height= .53cm, text width = 4.4cm, draw = gray!30]{\tikzdaytext};
    \draw (-1.8cm, -.1ex) node[anchor = west]{\footnotesize\pgfcalendarweekdayshortname{\pgfcalendarcurrentweekday}};
  },
  execute before day scope={
    \ifdate{day of month=1} {
      \pgftransformxshift{4.8cm}
      \draw (0,0)node [shape=rectangle, minimum height= .53cm, 
        text width = 4.4cm, fill = NavyBlue!80, text= white, draw = NavyBlue!80, text centered]
        {\textbf{\pgfcalendarmonthname{\pgfcalendarcurrentmonth}}};
    }{}
    \ifdate{workday}{\tikzset{every day/.style={fill=white}}}{}
    \ifdate{between=2026-07-21 and 2026-08-31}{%
        \tikzset{every day/.style={fill=gray!30}}}{}
    \ifdate{between=2026-11-02 and 2026-11-06}{%
        \tikzset{every day/.style={fill=gray!30}}}{}
    \ifdate{between=2026-12-21 and 2027-01-05}{%
        \tikzset{every day/.style={fill=gray!30}}}{}
    \ifdate{Saturday}{\tikzset{every day/.style={fill=blue!5}}}{}
    \ifdate{Sunday}{\tikzset{every day/.style={fill=red!10}}}{}
     \ifdate{equals=07-20}{\tikzset{every day/.style={fill=red!10}}}{}
    \ifdate{equals=08-11}{\tikzset{every day/.style={fill=red!10}}}{}
    \ifdate{equals=09-21}{\tikzset{every day/.style={fill=red!10}}}{}
    \ifdate{equals=09-22}{\tikzset{every day/.style={fill=red!10}}}{}
    \ifdate{equals=10-12}{\tikzset{every day/.style={fill=red!10}}}{}
    \ifdate{equals=11-03}{\tikzset{every day/.style={fill=red!10}}}{}
    \ifdate{equals=11-23}{\tikzset{every day/.style={fill=red!10}}}{}
  },
  execute at begin day scope={\pgftransformyshift{-.53*\pgfcalendarcurrentday cm}}
];
\termin{\year-07-20}{Marine Day}
\termin{\year-08-11}{Mountain Day}
\termin{\year-09-21}{Respect for the Aged Day}
\termin{\year-09-22}{Autumnal Equinox Day}
\termin{\year-10-12}{Sports Day}
\termin{\year-11-03}{Culture Day}
\termin{\year-11-23}{Labour Thanksgiving Day}
%%%%%%%%%%%%%%%%%%%
%予定を追加
\addschedule{20260717}{終業式}
\addschedule{20260901}{始業式}
\addschedule{20261218}{終業式}
\end{tikzpicture}
\end{center}
\end{document}
